\documentclass[UTF8,12pt,a4paper]{ctexart}
\usepackage[a4paper,margin=2.30cm,headheight=15pt]{geometry}
\usepackage{amsmath,amssymb,mathtools,bm}
\usepackage{booktabs,array,longtable}
\usepackage{xcolor,enumitem,fancyhdr,hyperref,listings}

\definecolor{sourcegray}{RGB}{75,84,97}
\definecolor{sourcebg}{RGB}{244,247,250}
\definecolor{linkblue}{RGB}{31,78,121}
\definecolor{keyblue}{RGB}{17,83,142}
\definecolor{keybg}{RGB}{239,247,255}
\hypersetup{colorlinks=true,linkcolor=linkblue,citecolor=linkblue,urlcolor=linkblue,
  pdftitle={第五题：投影法的有限体积离散与OpenFOAM求解器},pdfauthor={张云飞}}
\pagestyle{fancy}
\fancyhf{}
\lhead{第五题：不可压Navier--Stokes方程}
\rhead{投影法}
\cfoot{\thepage}
\setlength{\parindent}{2em}
\setlength{\parskip}{0.25em}
\setlist[itemize]{leftmargin=2.2em,itemsep=0.25em,topsep=0.35em}
\setlist[enumerate]{leftmargin=2.4em,itemsep=0.35em,topsep=0.35em}
\numberwithin{equation}{section}

\newcommand{\dd}{\,\mathrm d}
\newcommand{\dt}{\Delta t}
\newcommand{\Task}{\textnormal{[题目]}}
\newcommand{\Chorin}{\textnormal{[8]}}
\newcommand{\Guermond}{\textnormal{[9]}}
\newcommand{\Rone}{\textnormal{[R1]}}
\newcommand{\Rtwo}{\textnormal{[R2]}}
\newcommand{\Rthree}{\textnormal{[R3]}}
\newcommand{\source}[1]{\par\smallskip\noindent\colorbox{sourcebg}{%
  \parbox{0.955\linewidth}{\footnotesize\color{sourcegray}\textbf{参考来源：}#1}}\par\smallskip}
\newcommand{\keytitle}[1]{\par\smallskip\noindent\colorbox{keybg}{%
  \parbox{0.965\linewidth}{\color{keyblue}\textbf{重要公式：}#1}}\par\smallskip}
\lstdefinestyle{foam}{
  basicstyle=\ttfamily\small,frame=single,framerule=0.4pt,
  rulecolor=\color{sourcegray},backgroundcolor=\color{sourcebg},
  numbers=left,numberstyle=\tiny\color{sourcegray},numbersep=8pt,
  columns=fullflexible,keepspaces=true,showstringspaces=false,
  breaklines=true,xleftmargin=1.4em,framexleftmargin=1.0em}

\title{\bfseries 第五题：不可压Navier--Stokes方程\\投影法、有限体积离散与OpenFOAM求解器详细推导}
\author{张云飞}
\date{2026年8月27日}

\begin{document}
\maketitle
\begin{abstract}
本文严格围绕题目第5.1节，依据Chorin的原始分步思想\Chorin 与Guermond、Minev、Shen对投影方法的系统分类\Guermond，推导常密度不可压Navier--Stokes方程的一阶非增量压力校正投影法。随后按照Moukalled、Mangani、Darwish的有限体积离散主线\Rone，将预测动量方程、压力Poisson方程、单元速度修正和面体积通量修正逐项离散到一般非正交混合非结构网格，并给出与OpenFOAM算子、线性方程和时间循环对应的求解器骨架。全文沿用前四题统一记号：内部面全局方向为$O_f\to N_f$，$\bm S_f=|S_f|\bm n_f$，局部装配方向由$\sigma_{cf}$给出；不使用$d(f)$表示相邻单元，也不使用$A_f$表示面面积。本文只推导投影法；PISO另见独立文档。
\end{abstract}
\tableofcontents
\newpage

\section*{题目范围、参考资料和统一记号}
\addcontentsline{toc}{section}{题目范围、参考资料和统一记号}
\subsection*{题目要求}
题目第5.1节给出
\begin{align}
  \nabla\cdot\bm u&=0,\label{eq:continuity}\\
  \frac{\partial\bm u}{\partial t}+\nabla\cdot(\bm u\otimes\bm u)
  &=-\frac{1}{\rho}\nabla p+\frac{1}{\rho}\nabla\cdot(\mu\nabla\bm u).
  \label{eq:momentum}
\end{align}
题目要求分别基于投影法和PISO方法写出求解过程，再使用有限体积法及OpenFOAM提供的离散格式开发数值求解器。本文件完成其中的投影法部分。
\source{\Task，第5.1节；投影法指定参考文献为\Chorin、\Guermond。}

\subsection*{常密度与运动学压力}
题目算例中的$\rho$和$\mu$取常数。定义
\begin{equation}
  \boxed{\nu:=\frac{\mu}{\rho},\qquad \pi:=\frac{p}{\rho}.}
  \label{eq:kinematic}
\end{equation}
于是
\begin{align}
  \nabla\cdot\bm u&=0,\label{eq:continuity_pi}\\
  \frac{\partial\bm u}{\partial t}+\nabla\cdot(\bm u\otimes\bm u)
  &=-\nabla\pi+\nabla\cdot(\nu\nabla\bm u).
  \label{eq:momentum_pi}
\end{align}
OpenFOAM不可压求解器中的字段\texttt{p}通常存储$\pi=p/\rho$，因此后文数学公式用$\pi$，代码仍使用字段名\texttt{p}。若$\rho=1$，二者数值相同，但量纲含义仍应区分。
\source{\Rone，第15章不可压流动控制方程及压力--速度耦合；\Rtwo，不可压OpenFOAM求解器中运动学压力的量纲约定。}

\subsection*{统一记号}
\begin{longtable}{>{\centering\arraybackslash}p{0.21\linewidth}p{0.68\linewidth}}
\toprule
记号 & 含义\\
\midrule
$c$ & 当前控制体\\
$\Omega_c,V_c$ & 控制体$c$及其体积（二维为面积乘单位厚度）\\
$\mathcal F_c$ & 控制体$c$的全部面集合\\
$f$ & 内部面\\
$O_f,N_f$ & 面$f$的owner与neighbour单元\\
$\bm d_f=\bm x_{N_f}-\bm x_{O_f}$，$d_f=\lVert\bm d_f\rVert$ & 从$O_f$到$N_f$的中心连线及距离\\
$\bm n_f$ & 从$O_f$指向$N_f$的全局单位法向量\\
$\bm S_f=|S_f|\bm n_f$ & 从$O_f$指向$N_f$的全局有向面积向量\\
$\sigma_{cf}$ & $c=O_f$时为$+1$，$c=N_f$时为$-1$\\
$\bm S_{cf}=\sigma_{cf}\bm S_f$ & 面$f$相对于控制体$c$的外向面积向量\\
$F_f=\bm u_f\cdot\bm S_f$ & 按$O_f\to N_f$定向的体积通量\\
$Q_c^m=\sum_f\sigma_{cf}F_f+Q_c^{m,\mathrm{bnd}}$ & 未除体积的离散通量不平衡\\
$R_c^m=Q_c^m/V_c$ & 体积归一化的离散连续性残差\\
$\bm H_f^{\mathrm{adv}},\bm H_f^{\mathrm{diff}}$ & 全局定向的对流、黏性动量通量\\
\bottomrule
\end{longtable}
\begin{equation}
  \boxed{\sigma_{cf}=\begin{cases}+1,&c=O_f,\\-1,&c=N_f,\end{cases}
  \qquad\bm S_{cf}=\sigma_{cf}\bm S_f,\qquad\bm S_f=|S_f|\bm n_f.}
  \label{eq:orientation}
\end{equation}
\source{\Rtwo，第1.3节：有限体积网格、面积向量、owner--neighbour寻址和内部面反号装配；与前四题统一记号一致。}

\section{为什么不可压方程需要投影}
\subsection{压力约束与Helmholtz分解}
压力没有独立的演化方程，而是作为Lagrange乘子，使新速度落入无散空间
\[
  \mathcal V_0:=\{\bm v:\nabla\cdot\bm v=0\}.
\]
设预测速度分解为
\begin{equation}
  \bm u^*=\bm u^{n+1}+\dt\nabla\pi^{n+1}.
  \label{eq:helmholtz_scaled}
\end{equation}
于是
\begin{equation}
  \boxed{\bm u^{n+1}=\bm u^*-\dt\nabla\pi^{n+1}.}
  \label{eq:velocity_correction_cont}
\end{equation}
取散度并使用$\nabla\cdot\bm u^{n+1}=0$：
\begin{equation}
  \boxed{\nabla^2\pi^{n+1}=\frac{1}{\dt}\nabla\cdot\bm u^*.}
  \label{eq:pressure_poisson_cont}
\end{equation}
压力Poisson方程不是额外假设，而是由速度修正和无散约束直接推出。
\keytitle{投影法的核心闭环是式\eqref{eq:pressure_poisson_cont}与式\eqref{eq:velocity_correction_cont}。}
\source{\Chorin：中间速度、压力Poisson方程及速度修正；\Guermond，第2--3节：Helmholtz分解与非增量压力校正。}

\section{一阶非增量压力校正投影法}
\Guermond 将投影类方法分为压力校正、速度校正和一致分裂等类别。为与\Chorin 原始思想对应，本文选择一阶非增量压力校正，不混入增量式或旋转式。

\subsection{步骤一：预测速度}
\begin{equation}
  \boxed{\frac{\bm u^*-\bm u^n}{\dt}
  +\nabla\cdot(\bm u^n\otimes\bm u^*)
  -\nabla\cdot(\nu\nabla\bm u^*)=\bm0.}
  \label{eq:predictor_semiimplicit}
\end{equation}
对流输运速度取$\bm u^n$，被输运速度与黏性项取$\bm u^*$，形成线性半隐式预测方程。若全部取$n$层，则得到完全显式预测；投影步骤不变。半隐式形式便于用OpenFOAM的\texttt{fvm}算子构造稳定矩阵。
\source{\Chorin：无新压力的中间速度；\Guermond，第3节非增量压力校正；\Rone，第8、11、13章。半隐式线性化属于有限体积实现选择。}

\subsection{步骤二和三：压力与修正}
\begin{align}
  \boxed{\nabla^2\pi^{n+1}
  =\frac{1}{\dt}\nabla\cdot\bm u^*,}
  \label{eq:projection_poisson}\\
  \boxed{\bm u^{n+1}=\bm u^*-\dt\nabla\pi^{n+1}.}
  \label{eq:projection_corrector}
\end{align}
代回即有$\nabla\cdot\bm u^{n+1}=0$。

\subsection{压力边界和参考值}
若$\bm u^{n+1}=\bm g^{n+1}$，由速度修正式的法向分量得到
\begin{equation}
  \boxed{\frac{\partial\pi^{n+1}}{\partial n}
  =\frac{1}{\dt}(\bm u^*-\bm g^{n+1})\cdot\bm n.}
  \label{eq:pressure_neumann}
\end{equation}
不可穿透壁面且预测速度满足正确法向速度时，右端为零。纯Neumann问题只确定到常数，必须规定
\begin{equation}
  \boxed{\pi_{c_{\mathrm{ref}}}=0.}
  \label{eq:pressure_reference}
\end{equation}
\source{\Guermond：投影边界条件与分裂误差；\Rone，第15.6.2.5节：压力相对性。}

\section{预测动量方程的有限体积离散}
\subsection{控制体积分与面通量}
对式\eqref{eq:predictor_semiimplicit}在$\Omega_c$上积分，对散度项使用Gauss定理并采用面中心积分：
\begin{equation}
  \frac{V_c}{\dt}(\bm u_c^*-\bm u_c^n)
  +\sum_{f\in\mathcal F_c}(\bm u_f^n\cdot\bm S_{cf})\bm u_f^*
  -\sum_{f\in\mathcal F_c}\nu_f(\nabla\bm u^*)_f\cdot\bm S_{cf}
  =\bm0.
  \label{eq:predictor_fvm_local}
\end{equation}
用全局面方向写为
\begin{equation}
  \boxed{\frac{V_c}{\dt}(\bm u_c^*-\bm u_c^n)
  +\sum_{f\in\mathcal F_c^{\mathrm{int}}}\sigma_{cf}\bm H_f^{\mathrm{adv},*}
  -\sum_{f\in\mathcal F_c^{\mathrm{int}}}\sigma_{cf}\bm H_f^{\mathrm{diff},*}
  +\bm Q_c^{\mathrm{bnd},*}=\bm0,}
  \label{eq:predictor_fvm_global}
\end{equation}
其中
\begin{equation}
  \bm H_f^{\mathrm{adv},*}:=F_f^n\bm u_f^*,\quad
  F_f^n:=\bm u_f^n\cdot\bm S_f,\quad
  \bm H_f^{\mathrm{diff},*}:=\nu_f(\nabla\bm u^*)_f\cdot\bm S_f.
\end{equation}
\source{\Rone，第5章控制体积分和Gauss定理，第11章对流项，第8章扩散项，第13章非稳态项。}

\subsection{对流面值}
一阶迎风可写为
\begin{equation}
  \bm u_f^{\mathrm{up},*}=
  \begin{cases}
    \bm u_{O_f}^*,&F_f^n\ge0,\\
    \bm u_{N_f}^*,&F_f^n<0,
  \end{cases}
\end{equation}
或
\begin{equation}
  \boxed{\bm H_f^{\mathrm{adv},*}
  =(F_f^n)^+\bm u_{O_f}^*+(F_f^n)^-\bm u_{N_f}^*,}
  \label{eq:upwind_momentum_flux}
\end{equation}
其中$(F)^+=\max(F,0)$，$(F)^-=\min(F,0)$。若选\texttt{linearUpwind}，OpenFOAM用迎风值加梯度外推修正构造高阶面值，方向仍由$F_f$和$\sigma_{cf}$控制。
\source{\Rone，第11--12章；\Rtwo，OpenFOAM的\texttt{divSchemes}。}

\subsection{非正交扩散通量}
将
\begin{equation}
  \bm S_f=\bm E_f+\bm T_f,\qquad \bm E_f\parallel\bm d_f,
  \label{eq:nonorth_decomp}
\end{equation}
则
\begin{align}
  \bm H_f^{\mathrm{diff},*}
  &\approx D_f^u(\bm u_{N_f}^*-\bm u_{O_f}^*)
  +\nu_f(\nabla\bm u^*)_f\cdot\bm T_f,
  \label{eq:diff_flux_nonorth}\\
  D_f^u&:=\nu_f\frac{|\bm E_f|}{d_f}.
\end{align}
第一项隐式进入矩阵，第二项显式修正，对应\texttt{Gauss linear corrected}。
\source{\Rone，第8章非正交扩散；\Rtwo，\texttt{laplacianSchemes}和\texttt{snGradSchemes}。}

\subsection{代数方程}
每个速度分量形成
\begin{equation}
  \boxed{a_c^u\bm u_c^*+
  \sum_{f\in\mathcal F_c^{\mathrm{int}}}
  a_{cf}^u\bm u_{\mathrm{other}(c,f)}^*=\bm b_c^u.}
  \label{eq:predictor_algebra}
\end{equation}
“other”只用于文字说明，实际实现不定义$d(f)$：内部面直接以$O_f,N_f$成对装配。对角系数包含$V_c/\dt$，右端包含$(V_c/\dt)\bm u_c^n$及显式非正交修正。\texttt{fvm}算子自动完成矩阵装配。
\source{\Rone，第8、11、13章及第15.5.1节动量方程代数形式。}

\section{压力Poisson方程的有限体积离散}
\subsection{从离散连续性推出}
定义预测面通量$F_f^*:=\bm u_f^*\cdot\bm S_f$。为使离散连续性严格成立，直接修正面通量：
\begin{equation}
  \boxed{F_f^{n+1}
  =F_f^*-\dt(\nabla\pi^{n+1})_f\cdot\bm S_f.}
  \label{eq:flux_correction}
\end{equation}
对单元$c$要求
\begin{equation}
  Q_c^{m,n+1}:=
  \sum_{f\in\mathcal F_c^{\mathrm{int}}}\sigma_{cf}F_f^{n+1}
  +Q_c^{m,\mathrm{bnd},n+1}=0.
  \label{eq:discrete_continuity}
\end{equation}
代入得
\begin{equation}
  \boxed{\sum_{f\in\mathcal F_c}
  (\nabla\pi^{n+1})_f\cdot\bm S_{cf}
  =\frac{1}{\dt}\sum_{f\in\mathcal F_c}F_{cf}^*.}
  \label{eq:pressure_fvm}
\end{equation}
内部面$F_{cf}^*=\sigma_{cf}F_f^*$。这正是连续Poisson方程的守恒有限体积形式。
\keytitle{有限体积投影必须修正面通量，再逐控制体检查通量和。只修正单元速度后再插值到面，不能保证离散连续性严格成立。}
\source{\Chorin、\Guermond：Poisson方程由无散约束产生；\Rone，第15.5.2节：预测面通量代入连续性构造压力方程。}

\subsection{正交与非正交压力通量}
正交网格上
\begin{equation}
  (\nabla\pi)_f\cdot\bm S_f
  \approx\frac{|S_f|}{d_f}(\pi_{N_f}-\pi_{O_f}).
\end{equation}
非正交网格上
\begin{equation}
  (\nabla\pi)_f\cdot\bm S_f
  \approx\frac{|\bm E_f|}{d_f}(\pi_{N_f}-\pi_{O_f})
  +(\nabla\pi)_f\cdot\bm T_f.
  \label{eq:p_nonorth}
\end{equation}
第一项隐式入矩阵，第二项在非正交校正循环中显式更新。只有最后一次校正将压力方程通量用于更新$F_f^{n+1}$。
\source{\Rone，第8章和第15.5.2--15.5.3节；\Rtwo，\texttt{corrected}面法向梯度。}

\section{速度、通量修正与离散守恒}
\subsection{单元速度与面通量}
\begin{align}
  \boxed{\bm u_c^{n+1}
  =\bm u_c^*-\dt(\nabla\pi^{n+1})_c,}
  \label{eq:cell_velocity_correction}\\
  \boxed{F_f^{n+1}=F_f^*-H_f^p,\qquad
  H_f^p:=\dt(\nabla\pi^{n+1})_f\cdot\bm S_f.}
  \label{eq:final_phi}
\end{align}
单元梯度可用
\begin{equation}
  (\nabla\pi)_c\approx \frac{1}{V_c}
  \sum_{f\in\mathcal F_c}\pi_f\bm S_{cf}.
\end{equation}
$H_f^p$必须采用与压力方程相同的正交/非正交离散。

\subsection{离散守恒证明}
压力方程等价于
\[
  \sum_f\dt(\nabla\pi^{n+1})_f\cdot\bm S_{cf}
  =\sum_fF_{cf}^*.
\]
故
\begin{align}
  Q_c^{m,n+1}
  &=\sum_f\left[F_{cf}^*
  -\dt(\nabla\pi^{n+1})_f\cdot\bm S_{cf}\right]=0,\\
  \boxed{R_c^{m,n+1}:=\frac{Q_c^{m,n+1}}{V_c}=0.}
  \label{eq:discrete_zero_div}
\end{align}
实际“零”由压力线性求解器容差决定。

\section{边界条件}
方腔顶盖取$\bm u=(1,0)$，其余壁面$\bm u=\bm0$；三角腔顶边取$\bm u=(1,0)$，两条斜壁$\bm u=\bm0$。速度均为Dirichlet条件，不可穿透壁面满足$F_b=0$。

若预测步骤保持正确壁面法向速度，则壁面运动学压力取$\partial\pi/\partial n=0$，OpenFOAM中为\texttt{zeroGradient}。所有速度法向分量均指定时，压力矩阵有常数零空间，必须设置\texttt{pRefCell}和\texttt{pRefValue}。
\source{\Rone，第15.6节；\Guermond：投影法边界条件。}

\section{OpenFOAM求解器构造}
\subsection{公式与算子对应}
\begin{longtable}{p{0.37\linewidth}p{0.50\linewidth}}
\toprule
数学项 & OpenFOAM表达\\
\midrule
$\partial_t\bm u$ & \texttt{fvm::ddt(UStar)}\\
$\nabla\cdot(\bm u^n\otimes\bm u^*)$ & \texttt{fvm::div(phi, UStar)}\\
$\nabla\cdot(\nu\nabla\bm u^*)$ & \texttt{fvm::laplacian(nu, UStar)}\\
$\nabla\cdot\bm u^*$ & \texttt{fvc::div(phiStar)}\\
$\nabla\cdot(\dt\nabla\pi)$ & \texttt{fvm::laplacian(deltaT, p)}\\
$(\nabla\pi)_c$ & \texttt{fvc::grad(p)}\\
$\bm u_f\cdot\bm S_f$ & \texttt{fvc::flux(U)}\\
压力Laplacian面通量 & \texttt{pEqn.flux()}\\
\bottomrule
\end{longtable}

\subsection{核心代码骨架}
下列代码表达算法结构；字段读取、时间控制、包含文件和具体控制器接口应按OpenFOAM~14补齐。代码中的\texttt{p}表示$\pi$。
\begin{lstlisting}[style=foam,language=C++]
while (runTime.run())
{
    ++runTime;
    surfaceScalarField phiN(phi);

    // Step 1: pressure-free momentum predictor
    volVectorField UStar(U);
    fvVectorMatrix UStarEqn
    (
        fvm::ddt(UStar)
      + fvm::div(phiN, UStar)
      - fvm::laplacian(nu, UStar)
    );
    UStarEqn.solve();
    UStar.correctBoundaryConditions();

    surfaceScalarField phiStar(fvc::flux(UStar));

    // Step 2: projection pressure equation
    while (projection.correctNonOrthogonal())
    {
        fvScalarMatrix pEqn
        (
            fvm::laplacian(runTime.deltaT(), p)
         == fvc::div(phiStar)
        );
        pEqn.setReference(pRefCell, pRefValue);
        pEqn.solve();

        if (projection.finalNonOrthogonalIter())
        {
            phi = phiStar - pEqn.flux();
        }
    }

    // Step 3: cell-centred velocity correction
    U = UStar - runTime.deltaT()*fvc::grad(p);
    U.correctBoundaryConditions();

    // Keep conservative phi from the pressure correction.
    continuityErrs(phi);
    runTime.write();
}
\end{lstlisting}

\subsection{与PISO压力方程的区别}
PISO压力方程系数是动量矩阵对角倒数$r_A=1/a_c^u$，本投影法的修正系数是$\dt$。这里是
\[
  \nabla\cdot(\dt\nabla\pi)=\nabla\cdot\bm u^*,
\]
不能未经推导机械换成$\nabla\cdot(r_A\nabla\pi)=\nabla\cdot(\bm H/a)$。

\section{离散格式与线性求解配置}
\begin{lstlisting}[style=foam]
ddtSchemes
{
    default         Euler;
}
gradSchemes
{
    default         Gauss linear;
    grad(U)         cellLimited Gauss linear 1;
}
divSchemes
{
    default         none;
    div(phi,U)      bounded Gauss linearUpwind grad(U);
}
laplacianSchemes
{
    default         Gauss linear corrected;
}
interpolationSchemes
{
    default         linear;
}
snGradSchemes
{
    default         corrected;
}
fluxRequired
{
    default         no;
    p;
}
\end{lstlisting}

\begin{lstlisting}[style=foam]
solvers
{
    p
    {
        solver          GAMG;
        smoother        GaussSeidel;
        tolerance       1e-8;
        relTol          0.01;
    }
    UStar
    {
        solver          smoothSolver;
        smoother        symGaussSeidel;
        tolerance       1e-8;
        relTol          0.01;
    }
}
projection
{
    nNonOrthogonalCorrectors 1;
    pRefCell                 0;
    pRefValue                0;
}
\end{lstlisting}
容差和修正次数是算例配置，不是\Chorin 或\Guermond 规定的普适常数，应通过连续性残差和网格质量验证。
\source{\Rtwo、\Rthree：OpenFOAM字段、离散格式和求解字典；\Rone，第8、11、13、15章。}

\section{完整算法与检查}
给定$\bm u^0$、$\pi^0$和初始面通量，对$n=0,1,\ldots$执行：
\begin{enumerate}
  \item 根据$F_f^n$装配并求解预测动量方程\eqref{eq:predictor_fvm_global}，得到$\bm u^*$。
  \item 由$\bm u^*$构造预测面通量$F_f^*$。
  \item 装配压力Poisson方程\eqref{eq:pressure_fvm}，设置压力边界和参考压力。
  \item 在非正交循环内求解压力；只在最后一次使用压力矩阵通量更新式\eqref{eq:final_phi}。
  \item 使用式\eqref{eq:cell_velocity_correction}修正单元速度并重施边界条件。
  \item 计算$Q_c^{m,n+1}$和$R_c^{m,n+1}$，确认式\eqref{eq:discrete_zero_div}在容差内成立。
  \item 推进时间，直至速度场、主涡位置、残差和连续性误差均稳定。
\end{enumerate}

\keytitle{投影法有限体积求解器的四个核心式为：预测动量方程\eqref{eq:predictor_fvm_global}、压力方程\eqref{eq:pressure_fvm}、面通量修正式\eqref{eq:final_phi}和单元速度修正式\eqref{eq:cell_velocity_correction}。}

\begin{enumerate}
  \item 每个内部面只计算一次，owner加、neighbour减。
  \item 数学物理压力$p$与代码运动学压力$\pi=p/\rho$不得混淆。
  \item 压力方程修正面通量，不能再用简单插值覆盖守恒通量。
  \item 纯Neumann压力边界必须设置参考单元和值。
  \item 预测速度法向边界与压力Neumann条件必须满足式\eqref{eq:pressure_neumann}。
  \item 逐单元检查$R_c^m$，并检查全局边界净通量。
  \item 用中心线速度、主涡中心、流线和参考数据验证，不能只看残差。
\end{enumerate}

\section{参考来源与使用范围}
\begin{longtable}{p{0.13\linewidth}p{0.76\linewidth}}
\toprule
编号 & 本文实际使用内容\\
\midrule
\Task & 第5.1节控制方程和任务要求。\\
\Chorin & 无压力中间速度、压力Poisson方程、速度投影修正。\\
\Guermond & 投影方法分类；非增量压力校正；Helmholtz投影、边界和分裂误差。\\
\Rone & 有限体积积分、扩散、梯度、对流、时间、压力方程与面通量。\\
\Rtwo & OpenFOAM网格数据结构、owner--neighbour、面积向量、字段和离散格式。\\
\Rthree & OpenFOAM字典、运行和求解器设置。\\
\bottomrule
\end{longtable}

\begin{thebibliography}{99}
\bibitem{Chorin1967}
A.~J. Chorin,
“A numerical method for solving incompressible viscous flow problems,”
\emph{Journal of Computational Physics}, vol.~2, pp.~12--26, 1967.
\bibitem{Guermond2006}
J.-L. Guermond, P. Minev, and J. Shen,
“An overview of projection methods for incompressible flows,”
\emph{Computer Methods in Applied Mechanics and Engineering},
vol.~195, pp.~6011--6045, 2006.
\bibitem{Moukalled2016}
F. Moukalled, L. Mangani, and M. Darwish,
\emph{The Finite Volume Method in Computational Fluid Dynamics: An Advanced Introduction with OpenFOAM and Matlab},
Springer, 2016. 重点使用第5、8、9、11、13、15章。
\bibitem{Primer2021}
T. Holzmann and S. Rooney, \emph{The OpenFOAM Technology Primer}, 2021.
\bibitem{Holzinger2020}
G. Holzinger, \emph{OpenFOAM: A Little User-Manual}, 2020.
\end{thebibliography}
\end{document}
